\documentclass[a4paper,11pt]{article}
\usepackage[T1]{fontenc}
\usepackage[utf8]{inputenc}
\usepackage[left=2cm,right=2cm,top=2.5cm,bottom=2.5cm]{geometry}
\setlength{\headheight}{14pt}
\usepackage{xcolor}
\usepackage{mdframed}
\usepackage{parskip}
\usepackage{fancyhdr}
\usepackage{hyperref}

\definecolor{usercolor}{RGB}{30, 100, 200}
\definecolor{agentcolor}{RGB}{30, 150, 80}
\definecolor{doccolor}{RGB}{200, 90, 10}
\definecolor{userbg}{RGB}{235, 244, 255}
\definecolor{agentbg}{RGB}{235, 250, 238}
\definecolor{docbg}{RGB}{255, 243, 220}
\definecolor{answerbg}{RGB}{250, 250, 235}

\newmdenv[backgroundcolor=userbg,linecolor=usercolor,linewidth=1.5pt,
  leftmargin=0pt,rightmargin=80pt,innerleftmargin=10pt,innerrightmargin=10pt,
  innertopmargin=8pt,innerbottommargin=8pt,roundcorner=5pt,
  skipabove=6pt,skipbelow=3pt]{userbubble}

\newmdenv[backgroundcolor=agentbg,linecolor=agentcolor,linewidth=1.5pt,
  leftmargin=80pt,rightmargin=0pt,innerleftmargin=10pt,innerrightmargin=10pt,
  innertopmargin=8pt,innerbottommargin=8pt,roundcorner=5pt,
  skipabove=3pt,skipbelow=3pt]{agentbubble}

\newmdenv[backgroundcolor=docbg,linecolor=doccolor,linewidth=1pt,
  leftmargin=80pt,rightmargin=0pt,innerleftmargin=10pt,innerrightmargin=10pt,
  innertopmargin=6pt,innerbottommargin=6pt,roundcorner=4pt,
  skipabove=2pt,skipbelow=3pt]{docenv}

\newmdenv[backgroundcolor=answerbg,linecolor=gray,linewidth=0.5pt,
  leftmargin=80pt,rightmargin=0pt,innerleftmargin=10pt,innerrightmargin=10pt,
  innertopmargin=5pt,innerbottommargin=5pt,roundcorner=4pt,
  skipabove=2pt,skipbelow=6pt]{refanswerenv}

\pagestyle{fancy}
\fancyhf{}
\renewcommand{\headrulewidth}{0.3pt}
\cfoot{\thepage}

\lhead{\textcolor{gray}{\small RAG Conversation Dataset}}
\rhead{\textcolor{gray}{\small Conv.\,\#2\ --\ govt0611}}
\begin{document}

\begin{center}
{\LARGE\bfseries RAG Conversation #14}
\par\medskip
{\small\textcolor{gray}{Domain:\ \texttt{mt-rag-clapnq-elser-512-100-20240503}\quad|\quad Status:\ \textbf{accepted}}}
\end{center}
\vspace{0.4em}
\noindent\textcolor{gray}{\rule{\linewidth}{0.6pt}}
\vspace{0.4em}
\vspace{0.5em}
\noindent{\small\textcolor{usercolor}{\textbf{$\triangleright$\enspace User --- Turn~1}}}\hfill{\tiny\textcolor{gray}{Factoid \$\\textbackslash{}cdot\$ ANSWERABLE}}
\begin{userbubble}
name the viceroy at the time of quit india movement
\end{userbubble}
\noindent\hfill{\small\textcolor{agentcolor}{\textbf{Agent\enspace$\triangleleft$}}}
\begin{agentbubble}
The Viceroy at the time of the Quit India Movement was Lord Linlithgow. 
\end{agentbubble}
\noindent\hfill{\footnotesize\textcolor{doccolor}{\textbf{$\hookrightarrow$\enspace Retrieved documents~(3)}}}
\begin{docenv}
{\footnotesize\textcolor{doccolor}{\textbf{Doc~1}}}\hfill{\tiny\texttt{score:~10.317}}\ {\tiny\textcolor{gray}{\{'relevant': \{'f6ce971c-68d2-4c26-b1c0-477dd28a5181': \{'value': 'yes', 'timestamp': 1725377083\}\}\}}}\\[3pt]
{\small\textit{Quit India Movement}}\\[2pt]
\vspace{2pt}
{\footnotesize Quit India Movement
In 1939 , Indian nationalists were angry that British Governor - General of India , Lord Linlithgow , had brought India into the war without consultation with them . The Muslim League supported the war , but Congress was divided . Public}
\end{docenv}
\begin{docenv}
{\footnotesize\textcolor{doccolor}{\textbf{Doc~2}}}\hfill{\tiny\texttt{score:~8.121}}\ {\tiny\textcolor{gray}{\{'relevant': \{'f6ce971c-68d2-4c26-b1c0-477dd28a5181': \{'value': 'yes', 'timestamp': 1725377081\}, '87334bc6-dd39-4565-ae2b-832e60a24b08': \{'value': 'yes', 'timestamp': 1726869169\}\}\}}}\\[3pt]
{\small\textit{Partition of India}}\\[2pt]
\vspace{2pt}
{\footnotesize Partition of India
The Muslim League conducted its own investigation into the conditions of Muslims under Congress - governed provinces . The findings of such investigations increased fear among the Muslim masses of future Hindu domination . The view that Muslims would be unfairly treated in an independent India dominated by the Congress was now a part of the public discourse of Muslims . With the outbreak of World War II in 1939 , the viceroy , Lord Linlithgow , declared war on India 's behalf without consulting Indian leaders , leading the Congress provincial ministries to resign in protest . The Muslim League , which functioned under state patronage , in contrast , organized `` Deliverance Day '' , celebrations ( from Congress dominance ) and supported Britain in the war effort . When Linlithgow , met with nationalist leaders , he gave the same status to Jinnah as he did to Gandhi , and a month later described the Congress as a `` Hindu organization . ''}
\end{docenv}
\begin{docenv}
{\footnotesize\textcolor{doccolor}{\textbf{Doc~3}}}\hfill{\tiny\texttt{score:~7.691}}\ {\tiny\textcolor{gray}{\{'relevant': \{'f6ce971c-68d2-4c26-b1c0-477dd28a5181': \{'value': 'yes', 'timestamp': 1725377079\}, '87334bc6-dd39-4565-ae2b-832e60a24b08': \{'value': 'no', 'timestamp': 1726869148\}\}\}}}\\[3pt]
{\small\textit{Subhas Chandra Bose}}\\[2pt]
\vspace{2pt}
{\footnotesize Subhas Chandra Bose
On the outbreak of war , Bose advocated a campaign of mass civil disobedience to protest against Viceroy Lord Linlithgow 's decision to declare war on India 's behalf without consulting the Congress leadership . Having failed to persuade Gandhi of the necessity of this , Bose organised mass protests in Calcutta calling for the ' Holwell Monument ' commemorating the Black Hole of Calcutta , which then stood at the corner of Dalhousie Square , to be removed . He was thrown in jail by the British , but was released following a seven - day hunger strike . Bose 's house in Calcutta was kept under surveillance by the CID .}
\end{docenv}
\begin{refanswerenv}
{\footnotesize\textit{\textcolor{gray}{\textbf{Reference answer:}\enspace The Viceroy at the time of the Quit India Movement was Lord Linlithgow. He played a significant role in handling the situation during the Quit India Movement, including the arrest of much of the Indian National Congress leadership within hours of Gandhi's speech.}}}
\end{refanswerenv}
\vspace{0.5em}
\noindent{\small\textcolor{usercolor}{\textbf{$\triangleright$\enspace User --- Turn~2}}}\hfill{\tiny\textcolor{gray}{Summarization \$\\textbackslash{}cdot\$ Follow-up \$\\textbackslash{}cdot\$ ANSWERABLE}}
\begin{userbubble}
what did he do?
\end{userbubble}
\noindent\hfill{\small\textcolor{agentcolor}{\textbf{Agent\enspace$\triangleleft$}}}
\begin{agentbubble}
He had a significant role in the suppression of the Quit India Movement, which brought India into World War II without consultation with Indian nationalists, which led to their anger. The British had the support of his council. He brought India into World War II without consultation with Indian nationalists, which led to their anger.  The Conservatives in England were forced to concede some of the demands made by the Indians and On 8 August , the Viceroy issued a statement that has come to be referred as the `` August Offer '' . However , the Congress rejected the offer.
\end{agentbubble}
\noindent\hfill{\footnotesize\textcolor{doccolor}{\textbf{$\hookrightarrow$\enspace Retrieved documents~(3)}}}
\begin{docenv}
{\footnotesize\textcolor{doccolor}{\textbf{Doc~1}}}\hfill{\tiny\texttt{score:~21.061}}\ {\tiny\textcolor{gray}{\{'relevant': \{'f6ce971c-68d2-4c26-b1c0-477dd28a5181': \{'value': 'yes', 'timestamp': 1725377407\}\}\}}}\\[3pt]
{\small\textit{Quit India Movement}}\\[2pt]
\vspace{2pt}
{\footnotesize Quit India Movement
In 1939 , Indian nationalists were angry that British Governor - General of India , Lord Linlithgow , had brought India into the war without consultation with them . The Muslim League supported the war , but Congress was divided . Public}
\end{docenv}
\begin{docenv}
{\footnotesize\textcolor{doccolor}{\textbf{Doc~2}}}\hfill{\tiny\texttt{score:~18.140}}\ {\tiny\textcolor{gray}{\{'relevant': \{'f6ce971c-68d2-4c26-b1c0-477dd28a5181': \{'value': 'yes', 'timestamp': 1725377410\}\}\}}}\\[3pt]
{\small\textit{Quit India Movement}}\\[2pt]
\vspace{2pt}
{\footnotesize Quit India Movement
The Cripps Mission had failed , and on August 1942 , Gandhi made a call to Do or Die in his Quit India speech delivered in Bombay at the Gowalia Tank Maidan . The All - India Congress Committee launched a mass protest demanding what Gandhi called `` An Orderly British Withdrawal '' from India . Even though it was wartime , the British were prepared to act . Almost the entire leadership of the Indian National Congress was imprisoned without trial within hours of Gandhi 's speech . Most spent the rest of the war in prison and out of contact with the masses . The British had the support of the Viceroy 's Council ( which had a majority of Indians ) , of the All India Muslim League , the princely states , the Indian Imperial Police , the British Indian Army and the Indian Civil Service . Many Indian businessmen profiting from heavy wartime spending did not support the Quit India Movement . Many students paid more attention to Subhas Chandra Bose , who was in exile and supporting the Axis Powers . The only outside support came from the Americans , as President Franklin D. Roosevelt pressured Prime Minister Winston Churchill to give in to some of the Indian demands . The Quit India campaign was effectively crushed . The British refused to grant immediate independence , saying it could happen only after the war had ended .}
\end{docenv}
\begin{docenv}
{\footnotesize\textcolor{doccolor}{\textbf{Doc~3}}}\hfill{\tiny\texttt{score:~18.020}}\ {\tiny\textcolor{gray}{\{'relevant': \{'f6ce971c-68d2-4c26-b1c0-477dd28a5181': \{'value': 'yes', 'timestamp': 1725377413\}\}\}}}\\[3pt]
{\small\textit{Quit India Movement}}\\[2pt]
\vspace{2pt}
{\footnotesize Quit India Movement
In the meanwhile , crucial political events took place in England . Chamberlain was succeeded by Churchill as prime minister and the Conservatives , who assumed power in England , did not have a sympathetic stance towards the claims made by the Congress . In order to pacify the Indians in the circumstance of worsening war situation , the Conservatives were forced to concede some of the demands made by the Indians . On 8 August , the Viceroy issued a statement that has come to be referred as the `` August Offer '' . However , the Congress rejected the offer followed by the Muslim League .}
\end{docenv}
\begin{refanswerenv}
{\footnotesize\textit{\textcolor{gray}{\textbf{Reference answer:}\enspace Lord Linlithgow, the Viceroy of India during the Quit India Movement, brought India into World War II without consultation with Indian nationalists, which led to their anger. He also had a significant role in the suppression of the Quit India Movement, as almost the entire leadership of the Indian National Congress was imprisoned without trial within hours of Gandhi's speech launching the movement. The British had the support of the Viceroy's Council, the All India Muslim League, the princely states, the Indian Imperial Police, the British Indian Army, and the Indian Civil Service in their efforts to crush the movement.}}}
\end{refanswerenv}
\vspace{0.5em}
\noindent{\small\textcolor{usercolor}{\textbf{$\triangleright$\enspace User --- Turn~3}}}\hfill{\tiny\textcolor{gray}{Factoid, Keyword, Summarization \$\\textbackslash{}cdot\$ Follow-up \$\\textbackslash{}cdot\$ ANSWERABLE}}
\begin{userbubble}
opposition leader for the movement
\end{userbubble}
\noindent\hfill{\small\textcolor{agentcolor}{\textbf{Agent\enspace$\triangleleft$}}}
\begin{agentbubble}
A movement was launched at the Bombay session of the All - India Congress Committee by Mahatma Gandhi on 8 August 1942.  Many students paid attention to Subhas Chandra Bose as a leader,  who was in exile. He organized the Indian Legion in Germany , reorganized the Indian National Army with Japanese assistance , and soliciting help from the Axis Powers , conducted a guerrilla war against the British authorities.  
\end{agentbubble}
\noindent\hfill{\footnotesize\textcolor{doccolor}{\textbf{$\hookrightarrow$\enspace Retrieved documents~(3)}}}
\begin{docenv}
{\footnotesize\textcolor{doccolor}{\textbf{Doc~1}}}\hfill{\tiny\texttt{score:~26.073}}\ {\tiny\textcolor{gray}{\{'relevant': \{'f6ce971c-68d2-4c26-b1c0-477dd28a5181': \{'value': 'yes', 'timestamp': 1725378293\}\}\}}}\\[3pt]
{\small\textit{Quit India Movement}}\\[2pt]
\vspace{2pt}
{\footnotesize Quit India Movement
The Cripps Mission had failed , and on August 1942 , Gandhi made a call to Do or Die in his Quit India speech delivered in Bombay at the Gowalia Tank Maidan . The All - India Congress Committee launched a mass protest demanding what Gandhi called `` An Orderly British Withdrawal '' from India . Even though it was wartime , the British were prepared to act . Almost the entire leadership of the Indian National Congress was imprisoned without trial within hours of Gandhi 's speech . Most spent the rest of the war in prison and out of contact with the masses . The British had the support of the Viceroy 's Council ( which had a majority of Indians ) , of the All India Muslim League , the princely states , the Indian Imperial Police , the British Indian Army and the Indian Civil Service . Many Indian businessmen profiting from heavy wartime spending did not support the Quit India Movement . Many students paid more attention to Subhas Chandra Bose , who was in exile and supporting the Axis Powers . The only outside support came from the Americans , as President Franklin D. Roosevelt pressured Prime Minister Winston Churchill to give in to some of the Indian demands . The Quit India campaign was effectively crushed . The British refused to grant immediate independence , saying it could happen only after the war had ended .}
\end{docenv}
\begin{docenv}
{\footnotesize\textcolor{doccolor}{\textbf{Doc~2}}}\hfill{\tiny\texttt{score:~17.647}}\ {\tiny\textcolor{gray}{\{'relevant': \{'f6ce971c-68d2-4c26-b1c0-477dd28a5181': \{'value': 'yes', 'timestamp': 1725378122\}\}\}}}\\[3pt]
{\small\textit{Quit India Movement}}\\[2pt]
\vspace{2pt}
{\footnotesize Quit India Movement
After the onset of the war , only a group led by Subhas Chandra Bose took any decisive action . Bose organised the Indian Legion in Germany , reorganised the Indian National Army with Japanese assistance , and soliciting help from the Axis Powers , conducted a guerrilla war against the British authorities .}
\end{docenv}
\begin{docenv}
{\footnotesize\textcolor{doccolor}{\textbf{Doc~3}}}\hfill{\tiny\texttt{score:~14.933}}\ {\tiny\textcolor{gray}{\{'relevant': \{'f6ce971c-68d2-4c26-b1c0-477dd28a5181': \{'value': 'yes', 'timestamp': 1725380073\}\}\}}}\\[3pt]
{\small\textit{Quit India Movement}}\\[2pt]
\vspace{2pt}
{\footnotesize Quit India Movement
The Quit India Movement or the India August Movement , was a movement launched at the Bombay session of the All - India Congress Committee by Mahatma Gandhi on 8 August 1942 , during World War II , demanding an end to British Rule of India .}
\end{docenv}
\begin{refanswerenv}
{\footnotesize\textit{\textcolor{gray}{\textbf{Reference answer:}\enspace The key opposition to the Quit India Movement came from the Viceroy's Council, which had a majority of Indians, the All India Muslim League, the princely states, the Indian Imperial Police, the British Indian Army, the Indian Civil Service, and some Indian businessmen who were profiting from heavy wartime spending. Another opposition leader outside of these groups was Subhas Chandra Bose, who was in exile and supporting the Axis Powers. However, the only notable support for the movement came from the Americans, as President Franklin D. Roosevelt pressured Prime Minister Winston Churchill to give in to some of the Indian demands.}}}
\end{refanswerenv}
\vspace{0.5em}
\noindent{\small\textcolor{usercolor}{\textbf{$\triangleright$\enspace User --- Turn~4}}}\hfill{\tiny\textcolor{gray}{Summarization \$\\textbackslash{}cdot\$ Follow-up \$\\textbackslash{}cdot\$ ANSWERABLE}}
\begin{userbubble}
who is Patel?
\end{userbubble}
\noindent\hfill{\small\textcolor{agentcolor}{\textbf{Agent\enspace$\triangleleft$}}}
\begin{agentbubble}
Vallabhbhai Patel, popularly known as Sardar Patel , was the first Deputy Prime Minister of India . He was an Indian barrister and statesman , a senior leader of the Indian National Congress and a founding father of the Republic of India who played a leading role in the country 's struggle for independence and guided its integration into a united , independent nation . He acted as de facto Supreme Commander - in - chief of Indian army during the political integration of India and the Indo - Pakistani War of 1947 . He is credited by historians with ensuring the success of the rebellion across India. He was arrested and was imprisoned with the entire Congress Working Committee from 1942 to 1945 
\end{agentbubble}
\noindent\hfill{\footnotesize\textcolor{doccolor}{\textbf{$\hookrightarrow$\enspace Retrieved documents~(2)}}}
\begin{docenv}
{\footnotesize\textcolor{doccolor}{\textbf{Doc~1}}}\hfill{\tiny\texttt{score:~29.534}}\ {\tiny\textcolor{gray}{\{'relevant': \{'f6ce971c-68d2-4c26-b1c0-477dd28a5181': \{'value': 'yes', 'timestamp': 1725380645\}\}\}}}\\[3pt]
{\small\textit{Vallabhbhai Patel}}\\[2pt]
\vspace{2pt}
{\footnotesize Vallabhbhai Patel
Historians believe that Patel 's speech was instrumental in electrifying nationalists , who up to then had been skeptical of the proposed rebellion . Patel 's organising work in this period is credited by historians with ensuring the success of the rebellion across India . Patel was arrested on 9 August and was imprisoned with the entire Congress Working Committee from 1942 to 1945 at the fort in Ahmednagar . Here he spun cloth , played bridge , read a large number of books , took long walks , and practised gardening . He also provided emotional support to his colleagues while awaiting news and developments from the outside . Patel was deeply pained at the news of the deaths of Mahadev Desai and Kasturba Gandhi later that year . But Patel wrote in a letter to his daughter that he and his colleagues were experiencing `` fullest peace '' for having done `` their duty '' . Even though other political parties had opposed the struggle and the British had employed ruthless means of suppression , the Quit India movement was `` by far the most serious rebellion since that of 1857 '' , as the viceroy cabled to Winston Churchill . More than 100,000 people were arrested and many were killed in violent struggles with the police . Strikes , protests , and other revolutionary activities had broken out across India . When Patel was released on 15 June 1945 , he realised that the British were preparing proposals to transfer power to India .}
\end{docenv}
\begin{docenv}
{\footnotesize\textcolor{doccolor}{\textbf{Doc~2}}}\hfill{\tiny\texttt{score:~16.972}}\ {\tiny\textcolor{gray}{\{'relevant': \{'f6ce971c-68d2-4c26-b1c0-477dd28a5181': \{'value': 'yes', 'timestamp': 1725380641\}\}\}}}\\[3pt]
{\small\textit{Vallabhbhai Patel}}\\[2pt]
\vspace{2pt}
{\footnotesize Vallabhbhai Patel
Vallabhbhai Jhaverbhai Patel ( 31 October 1875 -- 15 December 1950 ) , popularly known as Sardar Patel , was the first Deputy Prime Minister of India . He was an Indian barrister and statesman , a senior leader of the Indian National Congress and a founding father of the Republic of India who played a leading role in the country 's struggle for independence and guided its integration into a united , independent nation . In India and elsewhere , he was often called Sardar , Chief in Hindi , Urdu , and Persian . He acted as de facto Supreme Commander - in - chief of Indian army during the political integration of India and the Indo - Pakistani War of 1947 .}
\end{docenv}
\begin{refanswerenv}
{\footnotesize\textit{\textcolor{gray}{\textbf{Reference answer:}\enspace Vallabhbhai Patel was a prominent leader in the Indian National Congress and a key figure in the Quit India Movement. He played a crucial role in electrifying nationalists and ensuring the success of the rebellion across India. After being arrested and imprisoned from 1942 to 1945, he provided emotional support to his colleagues and was deeply pained by the deaths of Mahadev Desai and Kasturba Gandhi. Upon his release, he realized that the British were preparing proposals to transfer power to India.}}}
\end{refanswerenv}
\vspace{0.5em}
\noindent{\small\textcolor{usercolor}{\textbf{$\triangleright$\enspace User --- Turn~5}}}\hfill{\tiny\textcolor{gray}{Factoid \$\\textbackslash{}cdot\$ Follow-up \$\\textbackslash{}cdot\$ UNANSWERABLE}}
\begin{userbubble}
how do you pronounce Vallabhbhai?
\end{userbubble}
\noindent\hfill{\small\textcolor{agentcolor}{\textbf{Agent\enspace$\triangleleft$}}}
\begin{agentbubble}
Sorry, but I cannot find the pronunciation of Vallabhbhai.
\end{agentbubble}
\noindent\hfill{\footnotesize\textcolor{doccolor}{\textbf{$\hookrightarrow$\enspace Retrieved documents~(2)}}}
\begin{docenv}
{\footnotesize\textcolor{doccolor}{\textbf{Doc~1}}}\hfill{\tiny\texttt{score:~37.053}}\ {\tiny\textcolor{gray}{\{'relevant': \{'f6ce971c-68d2-4c26-b1c0-477dd28a5181': \{'value': 'no', 'timestamp': 1725381526\}\}\}}}\\[3pt]
{\small\textit{Vallabhbhai Patel}}\\[2pt]
\vspace{2pt}
{\footnotesize Vallabhbhai Patel
Historians believe that Patel 's speech was instrumental in electrifying nationalists , who up to then had been skeptical of the proposed rebellion . Patel 's organising work in this period is credited by historians with ensuring the success of the rebellion across India . Patel was arrested on 9 August and was imprisoned with the entire Congress Working Committee from 1942 to 1945 at the fort in Ahmednagar . Here he spun cloth , played bridge , read a large number of books , took long walks , and practised gardening . He also provided emotional support to his colleagues while awaiting news and developments from the outside . Patel was deeply pained at the news of the deaths of Mahadev Desai and Kasturba Gandhi later that year . But Patel wrote in a letter to his daughter that he and his colleagues were experiencing `` fullest peace '' for having done `` their duty '' . Even though other political parties had opposed the struggle and the British had employed ruthless means of suppression , the Quit India movement was `` by far the most serious rebellion since that of 1857 '' , as the viceroy cabled to Winston Churchill . More than 100,000 people were arrested and many were killed in violent struggles with the police . Strikes , protests , and other revolutionary activities had broken out across India . When Patel was released on 15 June 1945 , he realised that the British were preparing proposals to transfer power to India .}
\end{docenv}
\begin{docenv}
{\footnotesize\textcolor{doccolor}{\textbf{Doc~2}}}\hfill{\tiny\texttt{score:~15.292}}\ {\tiny\textcolor{gray}{\{'relevant': \{'f6ce971c-68d2-4c26-b1c0-477dd28a5181': \{'value': 'no', 'timestamp': 1725381529\}\}\}}}\\[3pt]
{\small\textit{Mama and papa}}\\[2pt]
\vspace{2pt}
{\footnotesize Mama and papa
Though amma and appa are used in Tulu , they are not really Tulu words but used due to the influence of neighboring states ' languages . The actual words for mother in Tulu is appe pronounced IPA : ( ape ) and the word for father in Tulu is amme pronounced IPA : ( amæ ) . Note that the usage of these words is at odds with the usage pattern in other languages ( similar to Georgian in that sense ) . In Telugu , `` Thalli '' and `` Thandri '' are used for mother and father in formal Telugu . Amma and nana are used for mother and father for the informal way . Notice how nana refers to maternal grandfather in Hindi , and how that differs from its Telugu meaning . `` Nayana '' is also used for father in informal Telugu in the Rayalaseema region of Andhra Pradesh and Telangana of India . In Malayalam , `` Amma '' is used for mother . Hindus use `` Achan '' ( അച്ഛൻ ) for father and Christians use `` Appan '' . `` Appan '' was also previously used by matriarchal Hindu castes to denote the mother 's husband who may or may not be the biological father . The word `` Achan '' is a transformed Malayalam equivalent of the Sanskrit `` Arya '' for `` Sir / Master '' ( Arya - ajja - > ayya ) . `` Achan '' gained almost universal acceptance as the word for father among most Hindus during the 20th century . However , some families and communities also use ' appa ' etc. at least within their circles . The Malayali Christian communities often use a different set of words for father such as chaachchan ( ചാച്ചൻ ) , appachan ( അപ്പച്ചൻ ) etc . The Muslim communities use vaappa ( വാപ്പ ) , vaappachchi ( വാപ്പച്ചി ) , uppa ( ഉപ്പ ) etc . In Tamil , `` Thai '' and `` Thanthai '' are the formal Tamil words for mother and father ; informally Amma and Appa are much more common words for mother and father respectively . In the Kannada language , `` Thaai '' for mother and `` Thande '' for father are used formally . But to address them informally Kannadigas use amma or abbe for mother and appa for father .}
\end{docenv}
\begin{refanswerenv}
{\footnotesize\textit{\textcolor{gray}{\textbf{Reference answer:}\enspace The pronunciation of Vallabhbhai is approximately "Vuh-luhb-bahy." The first syllable "Vuh" is pronounced like the "u" in "up", the second syllable "luh" is pronounced like the "a" in "about", the third syllable "bh" is pronounced like the "b" in "boy" and the fourth syllable "bahy" is pronounced like the "i" in "it".}}}
\end{refanswerenv}
\vspace{0.5em}
\noindent{\small\textcolor{usercolor}{\textbf{$\triangleright$\enspace User --- Turn~6}}}\hfill{\tiny\textcolor{gray}{Summarization \$\\textbackslash{}cdot\$ Follow-up \$\\textbackslash{}cdot\$ ANSWERABLE}}
\begin{userbubble}
That is too bad.  Although the movement was non-violent, some ended up in violence, right? 
\end{userbubble}
\noindent\hfill{\small\textcolor{agentcolor}{\textbf{Agent\enspace$\triangleleft$}}}
\begin{agentbubble}
Yes, that's correct. While the Quit India Movement was primarily a non-violent movement, there were instances where it turned violent. Some protesters resorted to bombings, arson, cutting electricity and communication lines, and other destructive activities. They attacked 550 post offices , 250 railway stations , damaged many rail lines , destroyed 70 police stations , and burned or damaged 85 other government buildings.  The British response was also firm and over 100,000 arrests were made,  with mass detentions, fines, public flogging, and killings of hundreds of civilians.
\end{agentbubble}
\noindent\hfill{\footnotesize\textcolor{doccolor}{\textbf{$\hookrightarrow$\enspace Retrieved documents~(3)}}}
\begin{docenv}
{\footnotesize\textcolor{doccolor}{\textbf{Doc~1}}}\hfill{\tiny\texttt{score:~18.606}}\ {\tiny\textcolor{gray}{\{'relevant': \{'f6ce971c-68d2-4c26-b1c0-477dd28a5181': \{'value': 'yes', 'timestamp': 1725385575\}\}\}}}\\[3pt]
{\small\textit{Quit India Movement}}\\[2pt]
\vspace{2pt}
{\footnotesize Quit India Movement
attacked 550 post offices , 250 railway stations , damaged many rail lines , destroyed 70 police stations , and burned or damaged 85 other government buildings . There were about 2,500 instances of telegraph wires being cut . The greatest level of violence occurred in Bihar . The Government of India deployed 57 battalions of British troops to restore order .}
\end{docenv}
\begin{docenv}
{\footnotesize\textcolor{doccolor}{\textbf{Doc~2}}}\hfill{\tiny\texttt{score:~18.643}}\ {\tiny\textcolor{gray}{\{'relevant': \{'f6ce971c-68d2-4c26-b1c0-477dd28a5181': \{'value': 'yes', 'timestamp': 1725385579\}\}\}}}\\[3pt]
{\small\textit{Quit India Movement}}\\[2pt]
\vspace{2pt}
{\footnotesize Quit India Movement
The British swiftly responded with mass detentions . Over 100,000 arrests were made , mass fines were levied and demonstrators were subjected to public flogging . Hundreds of civilians were killed in violence many shot by the police army . Many national leaders went underground and continued their struggle by broadcasting messages over clandestine radio stations , distributing pamphlets and establishing parallel governments . The British sense of crisis was strong enough that a battleship was specifically set aside to take Gandhi and the Congress leaders out of India , possibly to South Africa or Yemen but ultimately did not take that step out of fear of intensifying the revolt .}
\end{docenv}
\begin{docenv}
{\footnotesize\textcolor{doccolor}{\textbf{Doc~3}}}\hfill{\tiny\texttt{score:~18.688}}\ {\tiny\textcolor{gray}{\{'relevant': \{'f6ce971c-68d2-4c26-b1c0-477dd28a5181': \{'value': 'yes', 'timestamp': 1725385582\}\}\}}}\\[3pt]
{\small\textit{Quit India Movement}}\\[2pt]
\vspace{2pt}
{\footnotesize Quit India Movement
One of the important achievements of the movement was keeping the Congress party united through all the trials and tribulations that followed . The British , already alarmed by the advance of the Japanese army to the India - Burma border , responded by imprisoning Gandhi . All the members of the Party 's Working Committee ( national leadership ) were imprisoned as well . Due to the arrest of major leaders , a young and until then relatively unknown Aruna Asaf Ali presided over the AICC session on 9 August and hoisted the flag ; later the Congress party was banned . These actions only created sympathy for the cause among the population . Despite lack of direct leadership , large protests and demonstrations were held all over the country . Workers remained absent en masse and strikes were called . Not all demonstrations were peaceful , at some places bombs exploded , government buildings were set on fire , electricity was cut and transport and communication lines were severed . Play}
\end{docenv}
\begin{refanswerenv}
{\footnotesize\textit{\textcolor{gray}{\textbf{Reference answer:}\enspace Yes, that's correct. While the Quit India Movement was primarily a non-violent movement, there were instances where it turned violent. Some protesters resorted to bombings, arson, cutting electricity and communication lines, and other destructive activities. The British response was also firm, with mass detentions, fines, public flogging, and killings of hundreds of civilians.}}}
\end{refanswerenv}
\vspace{0.5em}
\noindent{\small\textcolor{usercolor}{\textbf{$\triangleright$\enspace User --- Turn~7}}}\hfill{\tiny\textcolor{gray}{Factoid, Keyword \$\\textbackslash{}cdot\$ Follow-up \$\\textbackslash{}cdot\$ PARTIAL}}
\begin{userbubble}
population back then
\end{userbubble}
\noindent\hfill{\small\textcolor{agentcolor}{\textbf{Agent\enspace$\triangleleft$}}}
\begin{agentbubble}
I don't have the specific population figures for the time of the Quit India Movement, which took place from 1942 to 1945.  The Indian population rose to 300 million by 1920 and the population of undivided India in 1947 was approx 390 million .  The first post-colonial census , conducted in 1951 , counted 361.1 million people after partition in India. According to the 1951 Census of India , 2 \% of India 's population were refugees ( 1.3 \% from West Pakistan and 0.7 \% from East Pakistan ) . Delhi received the largest number of refugees for a single city -- the population of Delhi grew rapidly in 1947 from under 1 million ( 917,939 ) to a little less than 2 million ( 1,744,072 ) during the period 1941 -- 1951 
\end{agentbubble}
\noindent\hfill{\footnotesize\textcolor{doccolor}{\textbf{$\hookrightarrow$\enspace Retrieved documents~(4)}}}
\begin{docenv}
{\footnotesize\textcolor{doccolor}{\textbf{Doc~1}}}\hfill{\tiny\texttt{score:~18.837}}\ {\tiny\textcolor{gray}{\{'relevant': \{'f6ce971c-68d2-4c26-b1c0-477dd28a5181': \{'value': 'yes', 'timestamp': 1725385438\}\}\}}}\\[3pt]
{\small\textit{India}}\\[2pt]
\vspace{2pt}
{\footnotesize India
With 1,210,193,422 residents reported in the 2011 provisional census report , India is the world 's second-most populous country . Its population grew by 17.64 \% during 2001 -- 2011 , compared to 21.54 \% growth in the previous decade ( 1991 -- 2001 ) . The human sex ratio , according to the 2011 census , is 940 females per 1,000 males . The median age was 27.6 as of 2016 . The first post-colonial census , conducted in 1951 , counted 361.1 million people . Medical advances made in the last 50 years as well as increased agricultural productivity brought about by the `` Green Revolution '' have caused India 's population to grow rapidly . India continues to face several public health - related challenges .}
\end{docenv}
\begin{docenv}
{\footnotesize\textcolor{doccolor}{\textbf{Doc~2}}}\hfill{\tiny\texttt{score:~18.621}}\ {\tiny\textcolor{gray}{\{'relevant': \{'f6ce971c-68d2-4c26-b1c0-477dd28a5181': \{'value': 'yes', 'timestamp': 1725385528\}\}\}}}\\[3pt]
{\small\textit{Partition of India}}\\[2pt]
\vspace{2pt}
{\footnotesize Partition of India
According to the 1951 Census of India , 2 \% of India 's population were refugees ( 1.3 \% from West Pakistan and 0.7 \% from East Pakistan ) . Delhi received the largest number of refugees for a single city -- the population of Delhi grew rapidly in 1947 from under 1 million ( 917,939 ) to a little less than 2 million ( 1,744,072 ) during the period 1941 -- 1951 . The refugees were housed in various historical and military locations such as the Purana Qila , Red Fort , and military barracks in Kingsway Camp ( around the present Delhi University ) . The latter became the site of one of the largest refugee camps in northern India with more than 35,000 refugees at any given time besides Kurukshetra camp near Panipat . The camp sites were later converted into permanent housing through extensive building projects undertaken by the Government of India from 1948 onwards . A number of housing colonies in Delhi came up around this period like Lajpat Nagar , Rajinder Nagar , Nizamuddin East , Punjabi Bagh , Rehgar Pura , Jangpura and Kingsway Camp . A number of schemes such as the provision of education , employment opportunities , and easy loans to start businesses were provided for the refugees at the all - India level .}
\end{docenv}
\begin{docenv}
{\footnotesize\textcolor{doccolor}{\textbf{Doc~3}}}\hfill{\tiny\texttt{score:~18.576}}\ {\tiny\textcolor{gray}{\{'relevant': \{'f6ce971c-68d2-4c26-b1c0-477dd28a5181': \{'value': 'yes', 'timestamp': 1725385366\}\}\}}}\\[3pt]
{\small\textit{Partition of India}}\\[2pt]
\vspace{2pt}
{\footnotesize Partition of India
Massive population exchanges occurred between the two newly formed states in the months immediately following the Partition . `` The population of undivided India in 1947 was approx 390 million . After partition , there were 330 million people in India , 30 million in West Pakistan , and 30 million people in East Pakistan ( now Bangladesh ) . '' Once the lines were established , about 14.5 million people crossed the borders to what they hoped was the relative safety of religious majority . The 1951 Census of Pakistan identified the number of displaced persons in Pakistan at 7,226,600 , presumably all Muslims who had entered Pakistan from India . Similarly , the 1951 Census of India enumerated 7,295,870 displaced persons , apparently all Hindus and Sikhs who had moved to India from Pakistan immediately after the Partition . The two numbers add up to 14.5 million . Since both censuses were held about 3.6 years after the Partition , the enumeration included net population increase after the mass migration .}
\end{docenv}
\begin{docenv}
{\footnotesize\textcolor{doccolor}{\textbf{Doc~4}}}\hfill{\tiny\texttt{score:~16.676}}\ {\tiny\textcolor{gray}{\{'relevant': \{'f6ce971c-68d2-4c26-b1c0-477dd28a5181': \{'value': 'yes', 'timestamp': 1725385371\}\}\}}}\\[3pt]
{\small\textit{Famine}}\\[2pt]
\vspace{2pt}
{\footnotesize Famine
`` Population growth worsened the plight of the peasantry . As a result of peace and improved sanitation and health , the Indian population rose from perhaps 100 million in 1700 to 300 million by 1920 . While encouraging agricultural productivity , the British also provided economic incentives to have more children to help in the fields . Although a similar population increase occurred in Europe at the same time , the growing numbers could be absorbed by industrialization or emigration to the Americas and Australia . India enjoyed neither an industrial revolution nor an increase in food growing . Moreover , Indian landlords had a stake in the cash crop system and discouraged innovation . As a result , population numbers far outstripped the amount of available food and land , creating dire poverty and widespread hunger . '' -- - Craig A. Lockard , Societies , Networks , and Transitions}
\end{docenv}
\begin{refanswerenv}
{\footnotesize\textit{\textcolor{gray}{\textbf{Reference answer:}\enspace According to the document, the population of Karnataka in 1951 was 19,402,000. However, I don't have the specific population figures for the time of the Quit India Movement, which took place from 1942 to 1945.}}}
\end{refanswerenv}
\vspace{0.5em}
\noindent{\small\textcolor{usercolor}{\textbf{$\triangleright$\enspace User --- Turn~8}}}\hfill{\tiny\textcolor{gray}{Summarization \$\\textbackslash{}cdot\$ Follow-up \$\\textbackslash{}cdot\$ ANSWERABLE}}
\begin{userbubble}
when was the time Indian economy blossomed after independence?
\end{userbubble}
\noindent\hfill{\small\textcolor{agentcolor}{\textbf{Agent\enspace$\triangleleft$}}}
\begin{agentbubble}
The Indian economy started to blossom in the 1990s, following market-based economic reforms.  After 1991 economic liberalization, India achieved 6 - 7 \% average GDP growth annually . In FY 2015 and 2018 India 's economy became the world 's fastest growing major economy , surpassing China . 
\end{agentbubble}
\noindent\hfill{\footnotesize\textcolor{doccolor}{\textbf{$\hookrightarrow$\enspace Retrieved documents~(5)}}}
\begin{docenv}
{\footnotesize\textcolor{doccolor}{\textbf{Doc~1}}}\hfill{\tiny\texttt{score:~19.435}}\ {\tiny\textcolor{gray}{\{'relevant': \{'f6ce971c-68d2-4c26-b1c0-477dd28a5181': \{'value': 'yes', 'timestamp': 1725451078\}\}\}}}\\[3pt]
{\small\textit{Economy of India}}\\[2pt]
\vspace{2pt}
{\footnotesize Economy of India
The economy of India is a developing mixed economy . It is the world 's seventh - largest economy by nominal GDP and the third - largest by purchasing power parity ( PPP ) . The country ranks 139th in per capita GDP ( nominal ) with \$2,134 and 122nd in per capita GDP ( PPP ) with \$7,783 as of 2018 . After 1991 economic liberalisation , India achieved 6 - 7 \% average GDP growth annually . In FY 2015 and 2018 India 's economy became the world 's fastest growing major economy , surpassing China .}
\end{docenv}
\begin{docenv}
{\footnotesize\textcolor{doccolor}{\textbf{Doc~2}}}\hfill{\tiny\texttt{score:~18.898}}\ {\tiny\textcolor{gray}{\{'relevant': \{'f6ce971c-68d2-4c26-b1c0-477dd28a5181': \{'value': 'yes', 'timestamp': 1725451075\}, '87334bc6-dd39-4565-ae2b-832e60a24b08': \{'value': 'no', 'timestamp': 1726870583\}\}\}}}\\[3pt]
{\small\textit{Economy of India}}\\[2pt]
\vspace{2pt}
{\footnotesize Economy of India
The engineering industry of India includes its growing car , motorcycle and scooters industry , and productivity machinery such as tractors . India manufactured and assembled about 18 million passenger and utility vehicles in 2011 , of which 2.3 million were exported . India is the largest producer and the largest market for tractors , accounting for 29 \% of global tractor production in 2013 . India is the 12th - largest producer and 7th - largest consumer of machine tools .}
\end{docenv}
\begin{docenv}
{\footnotesize\textcolor{doccolor}{\textbf{Doc~3}}}\hfill{\tiny\texttt{score:~18.833}}\ {\tiny\textcolor{gray}{\{'relevant': \{'f6ce971c-68d2-4c26-b1c0-477dd28a5181': \{'value': 'yes', 'timestamp': 1725451073\}\}\}}}\\[3pt]
{\small\textit{India}}\\[2pt]
\vspace{2pt}
{\footnotesize India
In 2017 , the Indian economy was the world 's sixth largest by nominal GDP and third largest by purchasing power parity . Following market - based economic reforms in 1991 , India became one of the fastest - growing major economies and is considered a newly industrialised country . However , it continues to face the challenges of poverty , corruption , malnutrition , and inadequate public healthcare . A nuclear weapons state and regional power , it has the second largest standing army in the world and ranks fifth in military expenditure among nations . India is a federal republic governed under a parliamentary system and consists of 29 states and 7 union territories . It is a pluralistic , multilingual and multi-ethnic society and is also home to a diversity of wildlife in a variety of protected habitats .}
\end{docenv}
\begin{docenv}
{\footnotesize\textcolor{doccolor}{\textbf{Doc~4}}}\hfill{\tiny\texttt{score:~18.812}}\ {\tiny\textcolor{gray}{\{'relevant': \{'f6ce971c-68d2-4c26-b1c0-477dd28a5181': \{'value': 'yes', 'timestamp': 1725451070\}\}\}}}\\[3pt]
{\small\textit{Economy of India}}\\[2pt]
\vspace{2pt}
{\footnotesize Economy of India
India has one of the fastest growing service sectors in the world with an annual growth rate above 9 \% since 2001 , which contributed to 57 \% of GDP in 2012 -- 13 . India has become a major exporter of IT services , Business Process Outsourcing ( BPO ) services , and software services with \$154 billion revenue in FY 2017 . This is the fastest - growing part of the economy . The IT industry continues to be the largest private - sector employer in India . India is the third - largest start - up hub in the world with over 3,100 technology start - ups in 2014 -- 15 The agricultural sector is the largest employer in India 's economy but contributes to a declining share of its GDP ( 17 \% in 2013 -- 14 ) . India ranks second worldwide in farm output . The industry ( manufacturing ) sector has held a steady share of its economic contribution ( 26 \% of GDP in 2013 -- 14 ) . The Indian automobile industry is one of the largest in the world with an annual production of 21.48 million vehicles ( mostly two and three - wheelers ) in 2013 -- 14 . India had \$600 billion worth of retail market in 2015 and one of world 's fastest growing e-commerce markets .}
\end{docenv}
\begin{docenv}
{\footnotesize\textcolor{doccolor}{\textbf{Doc~5}}}\hfill{\tiny\texttt{score:~18.796}}\ {\tiny\textcolor{gray}{\{'relevant': \{'f6ce971c-68d2-4c26-b1c0-477dd28a5181': \{'value': 'yes', 'timestamp': 1725451068\}\}\}}}\\[3pt]
{\small\textit{Economy of India}}\\[2pt]
\vspace{2pt}
{\footnotesize Economy of India
The combination of protectionist , import - substitution , Fabian socialism , and social democratic - inspired policies governed India for sometime after the end of British rule . The economy was then characterised by extensive regulation , protectionism , public ownership of large monopolies , pervasive corruption and slow growth . Since 1991 , continuing economic liberalisation has moved the country towards a market - based economy . By 2008 , India had established itself as one of the world 's faster - growing economies .}
\end{docenv}
\begin{refanswerenv}
{\footnotesize\textit{\textcolor{gray}{\textbf{Reference answer:}\enspace The Indian economy started to blossom after independence in the 1990s, following market-based economic reforms. These reforms helped India become one of the fastest-growing major economies and transformed it into a newly industrialized country. However, it still faces challenges such as poverty, corruption, malnutrition, and inadequate public healthcare.}}}
\end{refanswerenv}
\end{document}
